\section{Worked Example 2: 同一模型，三种评估会给出三种结论}

假设一个模型有三个特征：
\begin{itemize}
    \item 知识题做得不错。
    \item 对开放式对话回答流畅。
    \item 价格比同类模型低。
\end{itemize}

如果你把它放在不同的评价框架里，结论会明显不同：

\begin{table}[H]
\centering
\begin{tabular}{p{3.3cm}p{4.2cm}p{5.4cm}}
\toprule
\textbf{评估方式} & \textbf{你会看到什么} & \textbf{可能结论} \\
\midrule
MMLU / GPQA & 知识覆盖与推理能力 & “能力不错” \\
Chatbot Arena & 人类偏好、回复风格 & “用户喜欢它” \\
Artificial Analysis & 能力-价格前沿 & “性价比高” \\
\bottomrule
\end{tabular}
\caption{同一模型在不同评估框架下会得到不同结论}
\end{table}

\begin{importantbox}{这不是矛盾，而是视角不同}
一个模型“在知识题上强”并不自动推出“用户体验最好”，更不自动推出“安全或最便宜”。评估本来就应该多维度并行。
\end{importantbox}

